\chapter{Lay Summary}
Data today often forms large networks, such as people connected through social media, money flowing between banks, vehicles moving across road systems.
These networks can grow enormous, sometimes holding billions of connections and reaching sizes needing terabytes or petabytes of storage.
Storing and analyzing data at that scale, quickly and without losing any of it, remains a difficult engineering problem.
Most organizations address it by stitching together several expensive, specialized software systems that often need many machines.

This dissertation introduces a single system that handles both storage and analysis together on one machine, even when the data is too large to fit in memory.
On standard performance tests, it matches or outperforms existing specialized systems that often require all data to fit in memory.
By replacing that expensive infrastructure with one tool on one machine, it cuts costs and makes it simpler to draw useful conclusions from large, connected datasets.